\PassOptionsToPackage{table}{xcolor}
\documentclass[pdflatex,iicol,sn-mathphys-ay]{sn-jnl}

\usepackage{graphicx}
\usepackage{amsmath}
\usepackage{amssymb}
\usepackage{comment}
\usepackage{algpseudocode}
\usepackage[linesnumbered,ruled,vlined]{algorithm2e}
\usepackage{multirow}
\usepackage[most]{tcolorbox}
\usepackage{float}
\usepackage{ulem}
\usepackage{xcolor}
\usepackage{subcaption}
\usepackage{stmaryrd}
\usepackage{tabularx}
\usepackage{booktabs}
\usepackage{makecell}
\usepackage{tikz}
\usepackage{array}
\usepackage{setspace}
\usepackage{pifont}

\newcommand{\cmark}{\ding{51}}
\newcommand{\xmark}{\ding{55}}
\newcommand{\omark}{\ding{108}}

\definecolor{darkblue}{rgb}{0.15,0.15,0.55}
\definecolor{lightgrey}{rgb}{0.75,0.75,0.75}
\newcommand{\qedm}{\hfill $\square$}
\DeclareUnicodeCharacter{25B3}{\ensuremath{\triangle}}
\newcommand\ny[1]{{\color{magenta}{#1}}}
\newcommand\sj[1]{{\color{violet}{#1}}}
\newcommand\sh[1]{{\color{orange}{#1}}}
\newcommand\gy[1]{{\color{purple}{#1}}}
\newcommand\cp[1]{{\color{blue}{#1}}}
\newcommand\jh[1]{{\color{green}{#1}}}

\hyphenation{op-tical net-works semi-conduc-tor}
\raggedbottom

\begin{document}

\title[BRL Architecture]{BRL Architecture}

\author*[1]{\fnm{Changmin} \sur{Park}}\email{cmp9877@gmail.com}
\author[1]{\fnm{Seongah} \sur{Park}}
\author[1]{\fnm{Nayoung} \sur{Kim}}
\author[1]{\fnm{Sungjin} \sur{Park}}
\author[1]{\fnm{Geonyeong} \sur{Ko}}
\author[2]{\fnm{Junho} \sur{Lee}}
\author[2]{\fnm{Changjoo} \sur{Nam}}

\affil*[1]{\orgname{Sogang University},
  \orgaddress{\city{Seoul}, \country{Republic of Korea}}}
\affil[2]{\orgname{Institution2},
  \orgaddress{\city{Seoul}, \country{Republic of Korea}}}

\abstract{We present a method for deciding when a robot should request external information and what it should ask during task execution. A robot operating with incomplete information may fail to determine the current state of the world, leading to task failure. However, resolving every such ambiguity through external queries would incur unnecessary costs. Therefore, a robot must determine both when to query and what information to request. In this paper, we maintain a weighted belief over possible task states during execution. The framework uses the uncertainty of this belief to decide when external information is needed and selects the hypothesis expected to reduce that uncertainty the most to decide what to ask. After receiving an answer from an external expert, the belief is updated and task execution continues. We evaluate the framework in two robotic task domains and compare it with alternative strategies for query timing and query selection. The results show that deciding when to query has a substantial effect on task success, whereas deciding what to query mainly affects the number of required queries. Compared with conformal-prediction-based and reward-based strategies, our method achieves higher task success without substantially increasing the number of queries.}

\keywords{Planning under uncertainty, online symbolic planning, robotic systems}

\maketitle

\input{section/01_intro}
\input{section/02_related}
\input{section/03_preliminaries}
\input{section/04_method}
\input{section/05_exp_setting}
\input{section/05_exp_system}
\input{section/06_conclusion}

\bmhead{Acknowledgements}
This research was supported in part by the National Research Foundation of Korea (NRF) grant funded by the Korea government (MSIT)(No. RS-2024-00461583) and the National Research Foundation of Korea (NRF) grant funded by the Korea government (MSIT) (No. RS-2024-00411007).

\bibliography{ref}

\begin{appendices}
\input{appendix/00_appx_content}
\input{appendix/01_algorithms.tex}
\input{appendix/02_detailedEXP.tex}
\end{appendices}

\end{document}
